% Point to the main file
\documentclass[../../Main.tex]{subfiles}

% ========== Start the document ==========
\begin{document}

\section{Section 2.4 - Sequences and Series}

\begin{definition}[Sequence, Explicit Formula / nth term]

    A \yellowbox{sequence} is an ordered list of numbers

    $$a_1, a_2, a_3, \ldots$$

    Equivalently, a sequence is a function $f: \setZ^+ \to \setR$

    We call $f$ the \yellowbox{explicit formula} or \yellowbox{nth term} of a sequence

\end{definition}

\begin{example}
    Determine the next 3 terms of the following sequences and the explicit formula

    \begin{enumerate}
        \item 
        {   
            $
                5\;\underbrace{\longrightarrow}_{+7}\;, 
                12\;\underbrace{\longrightarrow}_{+7}\;, 
                19\;\underbrace{\longrightarrow}_{+7}\;, 
                26\;\underbrace{\longrightarrow}_{+7}\;, 
                \ldots, 
                33\;\underbrace{\longrightarrow}_{+7}\;, 
                44\;\underbrace{\longrightarrow}_{+7}\;
            $
    
            \begin{solution}
                \begin{align*}
                    a_n &= 7n + b \\
                    5 = a_1 &= 7(1) + b \\
                    b &= 7 \\
                    \Aboxed{a_n &= 7n -2} \\
                \end{align*}
            \end{solution}
        }

        \item
        {
            $
                7\;\underbrace{\longrightarrow}_{\cdot 3}\;,
                21\;\underbrace{\longrightarrow}_{\cdot 3}\;,
                63\;\underbrace{\longrightarrow}_{\cdot 3}\;,
                189\;\underbrace{\longrightarrow}_{\cdot 3}\;,
                \ldots,
                567\;\underbrace{\longrightarrow}_{\cdot 3}\;,
                1701\;\underbrace{\longrightarrow}_{\cdot 3}\;,
                5103\;\underbrace{\longrightarrow}_{\cdot 3}\;
            $

            \begin{solution}
                \begin{align*}
                    a_n &= c3^n \\
                    a_1 &= 7 = c3^1 \\
                    c &= \frac{7}{3} \\
                    a_n &= \frac{7}{3} 3^n \\
                    \Aboxed{a_n &= 7 \cdot 3^{n-1}}
                \end{align*}
            \end{solution}
        }

        \item 
        {
            $
                1, 2, 6, 24, 120, 720, 5040
            $

            \boxed{a_n = n!}, read as "n factorial"

            $0! = 1$
        }

        \item 
        {
            $
                1,
                1,
                \;\underbrace{2}_{1 + 1},
                \;\underbrace{3}_{2 + 1},
                \;\underbrace{5}_{3 + 2},
                \;\underbrace{8}_{5 + 3},
                \;\underbrace{13}_{8 + 5},
                \;\underbrace{21}_{13 + 8},
            $

            \begin{solution}
                \begin{align*}
                    f_n = \frac{1}{\sqrt{5}} \left[ \left(\frac{1 + \sqrt{5}}{2} \right)^n - \left(\frac{1 - \sqrt{5}}{2} \right)^n\right] \tag{We'll derive this later}
                \end{align*}

                Another solution is

                \begin{align*}
                    a_1 &= 1 \\
                    a_2 &= 2 \\
                    a_n &= a_{n-1} + a_{n-2}, n >= 3 \\
                \end{align*}
            \end{solution}
        }
    \end{enumerate}

\end{example}

\begin{definition}
    A \textbf{recursive} sequence is defined in terms of the previous terms of the sequence
\end{definition}

\begin{example}

    Here are some examples of recursive sequences

    \begin{enumerate}
        \item 
        {
            $ a_1 = 2, a_2 = 3, a_n+1 = 5a_n - 6a_{n-1}, n >= 2 $

            \begin{align*}
                n = 2, a_3 &= 5a_2 - 6a_1 = 5 \cdot 3 - 6 \cdot 2 = 15 - 12 = 3 \\
                n = 3, a_4 &= 5a_3 - 6a_2 = 5 \cdot 3 - 6 \cdot 3 = 15 - 18 = -3 \\
                n = 4, a_5 &= 5a_4 - 6a_3 = 5 \cdot (-3) - 6 \cdot 3 = -15 - 18 = -33 \\
            \end{align*}
        }

        \item 
        {
            $ a_1 = 1, a_2 = 2, a_3 = 4, a_n = a_{n-1} + a_{n-2} - n - 3, n >= 4$
            
            \begin{align*}
                a_4 &= 4 + 2 - 1 = 5 \\
                a_5 &= 5 + 4 - 2 = 7 \\
                a_6 &= 7 + 5 - 4 = 8 \\
            \end{align*}
        }

    \end{enumerate}
\end{example}

\begin{definition}[Series]
    A \yellowbox{series} is the sum of an ordered list of numbers
    
    $$ a_1 + a_2 + a_3 + \ldots $$

    \[
    \sum_{k = 1}^{\infty} a_k
    \]

    Where:

    \begin{align*}
        k &= \text{Index} \\
        1 &= \text{Starting value} \\
        \infty &= \text{Ending value} \\
        a_k &= k^\text{th} \; \text{term} \\
    \end{align*}
    
\end{definition}

\begin{example}
    Find the sum of the series

    \item 
    {
        \[ 1 + 2 + 3 + \ldots + 100 \]

        \begin{solution}
            \[ 1 + 2 + 3 + \ldots + 98 + 99 + 100 \]
            \[ (1 + 100) + (2 + 99) + (3 + 98) + \ldots \]
            \[ (101) + (101) + (101) + \ldots \]

            Since there's 100 things we need to add, there's 50 pairs

            So 

            \begin{align*}
                \text{sum} &= 101(50) \\
                \Aboxed{ &= 5050 } \\
            \end{align*}
        \end{solution}
    }
\end{example}

\begin{example}

    But what about in general?
    
    $$1 + 2 + 3 + \ldots + n = ?$$

    We have two cases:
    
    \begin{enumerate}
        \item When $n$ is even
        \item When $n$ is odd
    \end{enumerate}

    If we assume n is even, then

    \begin{align}
        \Aboxed{1 + 2 + 3 + \ldots + n = (n + 1) \cdot \frac{n}{2}}
    \end{align}

    Since $(n + 1)$ gives the value of each pair and $\frac{n}{2}$ gives the number of pairs

    If we assume n is odd, then

    \begin{align*}
        \left[ 1 + 2 + 3 + \ldots + (n - 1) \right] + n &= n \cdot \frac{n - 1}{2} \\
        &= \frac{n^2 + n}{2} \\
        \Aboxed{&= (n + 1) \cdot \frac{n}{2}} \\
    \end{align*}

    since $n$ is the value of each pair $(n - 1 + 1)$ and $\frac{n - 1}{2}$ is the number of pairs

    \medskip 

    Another way to think about it is to let $S$ represent the sum

    \begin{align*}
        S &= 1 + 2 + 3 + \ldots + (n - 2) + (n - 1) + n \\
        S &= n + (n - 1) + (n - 2) + \ldots + 3 + 2 + 1 \tag{Switch Order} \\
        \hline \tag{Add} \\
        2S &= \underbrace{(n + 1) + (n + 1) + (n + 1) + \ldots + (n + 1) + (n + 1) + (n + 1)}_{n} \\
        \Aboxed{S &= (n + 1) \cdot \frac{n}{2}}
    \end{align*}

    Writing this is summation notation

    $$\sum_{k = 1}^{n} \frac{n(n + 1)}{2}$$

\end{example}

To Review...

\begin{enumerate}
    \item 
        \begin{align*}
            \sum_{k = 1}^{n} c = \underbrace{c + c + \ldots + c}_{n} = \Aboxed{nc} \tag{North Carolina Rule}
        \end{align*}

    \item 
        \begin{align*}
            \sum_{k = 1}^{n} k = 1 + 2 + 3 + \ldots + n = \Aboxed{(n + 1) \frac{n}{2}} \tag{Pairing Rule}
        \end{align*}

    \item 
        \begin{align*}
            \sum_{k = 1}^{n} k^2 = 1^2 + 2^2 + 3^2 + \ldots + n^2 = \Aboxed{\frac{n(n + 1)(2n + 1)}{6}} \tag{Prove Later}
        \end{align*}

    \item
        \begin{align*}
            \sum_{k = 1}^{n} k^3 = 1^3 + 2^3 + 3^3 + \ldots + n^3 = \Aboxed{\left[ \frac{n(n + 1)}{2} \right]^2} \tag{Prove Later}
        \end{align*}

\end{enumerate}

You \textbf{MUST} know these 4! They will be on exams!

\begin{example}
    Compute the Following

    \begin{enumerate}
        \item 
            \begin{align*}
                \sum_{k = 1}^{6} k = \frac{6}{2}(6 + 1) = \Aboxed{21} \tag{Use formula}
            \end{align*}

        \item 
            \begin{align*}
                \sum_{k = 1}^{6} 2k &= (2 \cdot 1) + (2 \cdot 2) + (2 \cdot 3) + (2 \cdot 4) + (2 \cdot 5) + (2 \cdot 6) \\
                &= 2(1 + 2 + 3 + 4 + 5 + 6) \\
                &= 2 \sum_{k = 1}^{6} k \\
                &= 2 \cdot 21 \\
                \Aboxed{&= 42} \\
            \end{align*}

        \begin{note}
            This means that we can \textbf{factor} in sums!
        \end{note}

        \item 
            \begin{align*}
                \sum_{k = 1}^{6} 2 = 6 \cdot 2 = \Aboxed{12} \tag{Use formula}
            \end{align*}

        \item 
            \begin{align*}
                \sum_{k = 1}^{6} (k + 2) &= (1 + 2) (2 + 2) + (3 + 2) + (4 + 2) + (5 + 2) + (6 + 2) \\
                &= (1 + 2 + 3 + 4 + 5 + 6) + (2 + 2 + 2 + 2 + 2 + 2) \\
                &= \sum_{k = 1}^{6} k + \sum_{k = 1}^{6} 2 \\
                &= 21 + 12 \\
                \Aboxed{&= 33} \\
            \end{align*}

        \begin{note}
            This means that in general,

            \begin{align*}
                \sum_{k = 1}^{n} ca_k &= c \sum_{k = 1}^{n} a_k \tag{Factor out constants} \\
                \sum_{k = 1}^{n} (a_k + b_k) &= \sum_{k = 1}^{n} a_k + \sum_{k = 1}^{n} b_k \tag{Split} \\
            \end{align*}
        \end{note}

        \item 
            \begin{align*}
                \sum_{k = 1}^{200} (k + 1)^2 &= \sum_{k = 1}^{200} (k^2 + 2k + 1) \tag{Expand} \\
                &= \sum_{k = 1}^{200} k^2 + 2 \sum_{k = 1}^{200} k + \sum_{k = 1}^{200} 1 \tag{Split up} \\
                &= \frac{200(201(401))}{6} + 2 \cdot \frac{200(201)}{2} + 200 \cdot 1 \tag{Use formulas} \\
                \Aboxed{&= 2727100} \tag{Use Calculator} \\ 
            \end{align*}

            We could also do

            \begin{align*}
                \sum_{k = 1}^{200} (k + 1)^2 &= 2^2 + 3^2 + 4^2 + \ldots 201^2 \\
                &= (1^2 + 2^2 + 3^2 + 4^2 + \ldots 201^2) - 1^2 \\
                &= \frac{201(202)(403)}{2} - 1 \\
                \Aboxed{&= 2727100} \tag{Use Calculator} \\ 
            \end{align*}

        \begin{theorem}[Constant Shift]
            This means that when you have a \textbf{constant shift} like $(k + 1)$ is a shift from $k$, then

            \begin{align*}
                \sum_{k = 1}^{n} f(k + a) = \sum_{k = 1}^{n + a} f(k) - \sum_{k = 1}^{a} f(k) \\
            \end{align*}

            \begin{center}
                or
            \end{center}

            \begin{align*}
                \sum_{k = m}^{n} a_k = \sum_{k = 1}^{n} a_k - \sum_{k = 1}^{m - 1} a_k 
            \end{align*}

            Where $f(k)$ is a function like $k^2$ that has it's own formula for the sum like

            \begin{align*}
                \sum_{k = 1}^{n} f(k) = \sum_{k = 1}^{n} k^2 = \frac{n(n + 1)(2n + 1)}{6} \\    
            \end{align*}

            Using this on the other example,

            \begin{align*}
                \sum_{k = 1}^{200} (k + 1)^2 &= \sum_{k = 1}^{201} k^2 - \sum_{k = 1}^{1} k^2  \\
                &= \frac{201(202)(403)}{2} - \frac{1(2)(1)}{2} \\
                &= 2727101 - 1 \\
                \Aboxed{&= 2727100} \tag{Use Calculator} \\ 
            \end{align*}

        \end{theorem}

        \item 
            \begin{align*}
                \sum_{k = 1}^{10} 2^k = S \\
                \\
                S &= 2^1 + 2^2 + 2^3 + \ldots + 2^9 + 2^{10} \\
                2S &= 2^2 + 2^3 + 2^4 + \ldots + 2^{10} + 2^{11} \\
                \hline \tag{Subtract} \\
                S - 2S &= 2^1 - 2^{11} \\
                S(1 - 2) &= 2^1 - 2^{11} \\
                S &= \frac{2 - 2^{11}}{1 - 2} = \frac{2 - 2048}{-1} \\
                \Aboxed{&= 2046} \\ 
            \end{align*}

        \begin{note}
            This means that

            \begin{align*}
                S = \sum_{k = 0}^{n} r^k \tag{When k starts at 0} \\
                \\
                S &= 1 + r^1 + r^2 + \ldots + r^{n - 1} + r^n \\
                rS &= 0 + r^1 + r^2 + r^3 + \ldots r^{n - 1} + r^n + r^{n + 1} \\
                \hline \tag{Subtract} \\
                S - rS &= 1 - r^{n + 1} \\
                S(1 - r) &= 1 - r^{n + 1} \\
                \Aboxed{S &= \frac{1 - r^{n + 1}}{1 - r}} \\
            \end{align*}

            \begin{align*}
                S = \sum_{k = 1}^{n} r^k \tag{When k starts at 1} \\
                \\
                S &= r^1 + r^2 + \ldots + r^{n - 1} + r^n \\
                &= r(1 + r^1 + r^2 + \ldots r^{n - 2} + r^{n - 1}) \\
                &= r \sum_{k = 1}^{n - 1}r^k \\
                &= r \cdot \frac{1 - r^{n - 1 + 1}}{1 - r} \\ 
                \Aboxed{&= r \cdot \frac{1 - r^n}{1 - r}} \\
            \end{align*}

            In general,

            \begin{align*}
                \Aboxed{\sum_{k = \text{start}}^{n} r^k = \frac{r^{\text{start}} - r^{n + 1}}{1 - r}}
            \end{align*}
            
            \begin{note}
                This works for $r \in \setR, (-\infty, \infty)$, including \textbf{negative} and \textbf{fractions!}

                This can also work for sums going to $\infty$, remember that a fraction to $\infty$ is just $0$
            \end{note}
        \end{note}

        There's more examples, but I'm too lazy to write them, so I'll just write the equations

        \begin{note}
            \begin{align*}
                \sum_{k = 0}^{\infty} kr^k = \sum_{k = 1}^{\infty} kr^k = \; \Aboxed{\frac{r}{(1 - r)^2}} \\
            \end{align*}
        \end{note}

    \end{enumerate}
\end{example}

% ========== End the document ==========
\end{document}